% Generated draft: requires mathematical and source review.
\documentclass[11pt]{article}
\usepackage[margin=25mm]{geometry}
\usepackage{fontspec}
\setmainfont{Latin Modern Roman}
\usepackage{amsmath,amssymb,mathtools,unicode-math}
\setmathfont{Latin Modern Math}
\usepackage{enumitem,needspace,xurl,hyperref}
\setlength{\parindent}{0pt}
\setlength{\parskip}{5pt}
\setlength{\emergencystretch}{3em}
\newcommand{\N}{\mathbb{N}}
\newcommand{\Z}{\mathbb{Z}}
\newcommand{\Q}{\mathbb{Q}}
\newcommand{\R}{\mathbb{R}}
\newcommand{\C}{\mathbb{C}}
\newcommand{\F}{\mathbb{F}}
\newcommand{\A}{\mathbb{A}}
\newcommand{\PP}{\mathbb{P}}
\newcommand{\E}{\mathbb{E}}
\newcommand{\eps}{\varepsilon}
\newcommand{\dd}{\,\mathrm{d}}
\begin{document}
\section*{UnsolvedMath ID: 3800020}
\textbf{Source problem code:} catalogue entry
\subsection{Definitions and mathematical statement}
A convex polytope realization specifies a convex polytope in Euclidean space with the prescribed metric or combinatorial data. The source groups five constructive problems, customarily concerning three-dimensional polytopes:
\begin{itemize}
\item Aleksandrov: reconstruct from a planar net or intrinsically glued polygonal surface, under the intended gluing and curvature hypotheses.
\item Cauchy: reconstruct from facet shapes together with their adjacency/gluing data, in the compatible convex setting.
\item Minkowski: reconstruct from outward normal directions and the corresponding facet areas. In a standard full-dimensional three-dimensional form, positive areas $a_i$ and unit normals $u_i$ satisfy $\sum_i a_i u_i=0$ and span $\R^3$; the polytope is determined up to translation.
\item Steinitz: realize a $3$-connected planar graph as the edge graph of a convex $3$-polytope.
\item Koebe: realize the intended circle-packing or sphere-tangency refinement associated with a planar graph. The precise polytope-tangency version is not stated in the excerpt.
\end{itemize}
These descriptions identify the separate input types; the omitted theorem hypotheses and degeneracy conventions must be recovered rather than inferred to be interchangeable. Metric reconstruction may have uniqueness up to rigid motion or translation, whereas graph realization alone does not prescribe a unique shape.

\textbf{Algorithmic problem.} Develop exact polynomial-time algorithms for the constructive forms of these five named realization theorems, reconstructing a convex polytope from the applicable net, facets with adjacency, area-weighted normals, or $3$-connected planar edge graph. Each input must be interpreted with the compatibility assumptions of its theorem; the record does not separately ask for feasibility testing of arbitrary inconsistent data.

The term exact requires an output representation: rational coordinates, algebraic numbers, or another exact encoding need not have the same computational cost. Polynomial time likewise needs an input encoding and cost model. These are not supplied, nor are geometric nondegeneracy and allowable degeneracies. A small numerical coordinate bound differs from a polynomial coordinate bit-length bound, and neither is explicitly required in the literal statement.

The supplied summary reports constructive methods via convex optimization for Minkowski, circle packing for Koebe, and a polynomial-time high-precision Steinitz realization, but says uniform exact/small-coordinate guarantees are not established for the whole family, particularly naming Steinitz, Cauchy, and Aleksandrov. No algorithms, precise accuracy models, or citations are supplied. Those scope claims require checking individually and are not taken as verified lower bounds or impossibility results. The background warns that the source pages have not been systematically updated since 2001.
\subsection{Short English statement}
Give exact polynomial-time constructions for five classical convex-polytope realization theorems, using their distinct metric or graph inputs. The exact output encoding, cost model, compatibility assumptions, degeneracies, and Koebe variant are unspecified, so the scope of reported partial algorithms needs review.
\subsection{Sources}
\begin{itemize}
\item Local UnsolvedMath registry, permanent ID 3800020, source code catalogue entry, ``Building convex polytopes''; attribution field: Jeff Erickson and cited proposers. Supplied status partially solved, classification PARTIAL-PROGRESS, updated 2026-08-06.
\item Erickson, \emph{Open Problems in Computational Geometry}, item ``Building convex polytopes,'' supplied index \url{https://jeffe.cs.illinois.edu/open/}. Reported HTML extraction and access: 2026-07-29. The source explicitly warns that its pages have not been systematically updated since 2001. Rights: NEEDS\_REVIEW; public access verified but redistribution permission not established.
\item The statement names Aleksandrov, Cauchy, Minkowski, Steinitz, and Koebe by their realization theorems. The summary's algorithmic progress has no complete primary references or precise complexity/precision statements in the supplied record.
\end{itemize}
\end{document}
